\section{Introduction}
An evaluator observes that images prompted ``in the style of Monet'' move
closer to Monet's reference collection. What should that observation establish?
A painting instruction can improve proximity to several related painters at
once. A classifier can recover prompted names from differences peculiar to a
generator. A high-confidence reference assignment can still identify the wrong
prompt. These possibilities require different evaluation checks.

Our main audit crosses six requested configurations with four related painters,
generic-painting and artist-free controls, 14 fixed scenes and two repeats
(1,008 images). A separate retained collection adds 2,000 SD-Turbo images under
a different template and 25 matched-seed blocks. We also test whether a rule
calibrated on historical artworks can withhold unreliable prompt-name assignments.
The targets are digital reference measurements and recorded prompt names.
Historical content and reproduction processes enter the finite reference target;
neither target establishes perceived resemblance. The additional analyses reuse
previously exposed data, with their new rules fixed before their outcomes.

The empirical question is how much an evaluation conclusion changes when its
target is made explicit. We separate common naming movement from labeled painter
contrasts, then separate reference-based decisions from within-generator label
structure. Common naming movement and a generated-to-reference mixture offset
are distinct: the first explains proximity gain over a prompt control, while
the second can change class intercepts without changing painter contrasts.
Table~\ref{tab:quantity-map} maps the required checks. These use established
geometry and decision rules; the contribution is the measured evaluation
behavior and its limits, not a new algorithm or perceptual metric.

\begin{table*}[t]
\caption{Evaluation questions and the checks needed to interpret their readouts.
All targets concern recorded images or prompt names, not perceptual fidelity.}
\label{tab:quantity-map}
\centering\small
\begin{tabularx}{\linewidth}{@{}p{.25\linewidth}p{.25\linewidth}X@{}}
\toprule
Question & Readout & Interpretation check\\\midrule
Does naming increase reference proximity? & Own-painter prototype gain & Separate common movement from labeled agreement.\\
Do painter differences match the reference? & Aligned amplitude $\beta$ and error $D$ & State the feature units, reference target and scene aggregation.\\
Can the prompt name be recovered? & Top-1 accuracy & Compare reference decisions, fixed mean translation and supervised generated centroids.\\
When should a name assignment be issued? & Accepted error and coverage & Test a fixed reference-calibrated gate against margin filtering at equal coverage.\\
\bottomrule
\end{tabularx}
\end{table*}


\section{Related Work}
\citet{somepalli2024style} introduce CSD and General Style Similarity: a generated
embedding's dot product with an averaged artist prototype, including
content-constrained prompts. \citet{su2025prompted} evaluate prompted-artist
recognition using content-controlled name substitutions and matched artist-free
images; generated-image retrieval usually outperforms real-image retrieval in
their benchmark. Our stronger generated-prototype comparator therefore
contextualizes an established domain distinction. We audit common and labeled
contributions to proximity gain, not a new recognition task.

\citet{frochte2026csd} diagnose raw CSD cosine and shared-tradition confusions,
then evaluate local-density readouts and input-resolution changes. Their
absolute-score and pair-verification study excludes closed-set discrimination.
Our controlled name interventions quantify common/labeled gain and connect a
separate mean offset to prompted-name decisions. The general warning that raw
style similarity needs validation is established, not our novelty claim.

\citet{deliege2025} assess stylistic imitation and stereotyping with experts.
AI-Pastiche studies authenticity and prompt adherence \citep{asperti2025pastiche};
AI-WikiArt uses painting-derived captions and vision-language models for
attribution \citep{fu2025aiwikiart}. Our measured painter contrasts and same-image
CLIP/CSD comparisons target neither convincing imitation nor physical authorship.

Quantitative art research examines color, composition and image statistics
\citep{kim2014,sigaki2018,lee2020}. Unlike the contextual image-to-image
experiments of \citet{kim2026}, our text-only intervention supplies no original
composition.

Generative evaluation separates fidelity and diversity \citep{naeem2020},
while conditional metrics distinguish within-condition variation from relationships
among means \citep{benny2021conditional}. We use independent-partition inner products
as the primary noise-corrected squared-difference estimator \citep{diedrichsen2017representational}
and energy statistics as a complementary distributional measure \citep{szekely2013}.
Processing sensitivity \citep{parmar2022} and limits of learned-embedding
separation \citep{asperti2026} motivate the source-quality checks and the
restriction to a specified digital reference target.

Target-domain centroid classifiers also address classifier bias under domain
shift \citep{liang2021atdoc}. Our retrospective decision diagnostic uses a fixed
mean translation and supervised centroids with frozen features; it is not an
adaptation-method contribution.
Selective classification trades prediction coverage for error
\citep{geifman2017selective}. We test a fixed reference-calibrated rejection rule;
distribution-shift guarantees require assumptions beyond those established for
these historical and generated panels \citep{tibshirani2019shift}.

\section{Design and Measurements}\label{sec:design}
\subsection{Painters and Reference Images}
The primary panel covers four related painters through 649 distinct reproductions: 297 Monet,
106 Sisley, 141 Pissarro and 105 C\'ezanne. Artist means receive equal weight,
so the larger Monet collection does not dominate the between-artist target.
A separate 221-work development panel sets common feature units; no primary
reference or generated image is used to fit that transformation. Both panels
were assembled before the present experiment and had been used in earlier
analyses. They are explicit, finite comparison panels rather than newly withheld
samples of each painter's complete oeuvre.

Reference files carry Wikimedia Commons source metadata and Wikidata work
identities \citep{commonsglam,vrandecic2014wikidata}. These support traceable
selection and rights records, but do not standardize photography, color
calibration or conservation state. These are measurements of digital
reproductions, whose capture conditions can affect the comparison.

\subsection{A Common-scene, Six-model Experiment}
We compare six requested service configurations, labeled GPT Image 1, GPT Image 2, GPT Image 2.5 Flare (Flare),
GPT Image 2.5 Sunburst (Sunburst), Nano Banana 2 and FLUX.2 Max (FLUX).
Fourteen fixed outdoor scene descriptions (briefs) span water,
built, land and mixed compositions. Each scene is crossed with six clauses:
no painting instruction; ``Render as an oil painting''; and ``Render as an oil
painting in the style of [painter]'' for each of the four painters. The scene
text is otherwise identical. Each model--scene--clause cell receives two
separately requested outputs, giving 1,008 images collected on 10 September 2026 (UTC).

All models receive text only and return 1,024-by-1,024-pixel images.
The four OpenAI requests specify medium quality; Nano Banana 2 specifies
1K resolution, while FLUX requests PNG output without an explicit resolution
field. These comparisons concern the recorded
configurations, not maximum attainable quality or equal expenditure per image.
Response records do not echo model identifiers, so these labels identify requested
configurations rather than independently verified checkpoints
(Appendix~\ref{app:provenance}).
The exact prompts, parameters and randomized request order were fixed before
collection. The 14 scenes are a fixed design,
not a random sample of all possible content.
Earlier, separate prompt interventions are summarized in Appendix~\ref{app:cohorts}.

Appendix Figure~\ref{fig:original-generated} shows all six conditions for one brief.
GPT Image 2's C\'ezanne trees have darker blue-green outlines than its Monet
and Sisley trees: a visible name response in this scene, whose agreement with
historical painter differences remains to be measured.



\subsection{What the Features Measure}\label{sec:features}
The 31 features in Table~\ref{tab:features} capture properties that need not
change together. Two images can have similar chroma (color strength) but different
hue diversity, or similar total contrast but different edge directions and
fine-scale texture. Such distinctions make results by feature family interpretable.
There are 11 color coordinates (lightness, chroma, hue and color increments),
eight spatial coordinates (spectra, gradients, orientations and quadrants),
and 12 texture coordinates (wavelet energy, local binary patterns and local contrast).
Correlated coordinates and sensitivity to image processing nevertheless prevent
treating Euclidean distance as a validated measure of artistic style.



Images undergo orientation correction and conversion to sRGB using valid
embedded profiles; compatible unprofiled files are treated as sRGB. The primary
pipeline preserves aspect ratio at a 512-pixel short side, using Lanczos resizing
without upsampling. No painting-region mask is applied: source margins or
calibration strips can enter the measurements. Each coordinate is centered and scaled by its
median and interquartile range in the development panel, with equal total
weight per painter. Thus a unit is a development-panel IQR in that coordinate,
not a perceptual just-noticeable difference. The Euclidean metric weights
standardized coordinates equally, without further equalizing feature families.



\section{Comparing Painter Differences}\label{sec:method}
\subsection{What Counts as Agreement?}\label{sec:contrast-intuition}
Subtract each set's four-painter mean to isolate labeled differences. For
example, reference Monet is $.325$ development-IQR units darker than Sisley
in median lightness; GPT Image 2's mean gap is $.254$ in that direction.
A shared brightness shift cannot create this contrast.
Aligned amplitude $\beta$ measures agreement along the reference pattern;
error $D$ also counts differences in other directions. No distinction, exact
agreement and doubled reference contrasts have expected errors 1, 0 and 1.
Thus stronger responses need not give lower error.

\subsection{Computing the Two Scores}
Let $\mu_a\in\R^{31}$ be painter $a$'s standardized reference mean and
$r_a=\mu_a-\frac14\sum_b\mu_b$ its contrast. The total squared size of these
contrasts, $H=\sum_a\|r_a\|^2$, normalizes both scores. Here $\|u\|^2$ sums squared
coordinates, $\langle u,v\rangle$ sums matching-coordinate products, and
hats mark estimates. For generated measurements $z_{msak}$, indices denote
model $m$, scene $s$, painter $a$ and repeat $k\in\{1,2\}$:
\begin{equation}
 d_{msak}=z_{msak}-\frac14\sum_b z_{msbk}.
 \label{eq:center}
\end{equation}
Centering removes shared additive shifts, but shared rescaling or mixing
of coordinates can still change the contrasts
(Appendix~\ref{app:score-details}). With $S=14$ scenes, the aligned response is
\begin{equation}
\begin{gathered}
\widehat\beta_{msk}=\frac{\sum_a\langle d_{msak},r_a\rangle}{H},\\
\widehat\beta_m=\frac1S\sum_s\frac{\widehat\beta_{ms1}+\widehat\beta_{ms2}}2.
\end{gathered}
\label{eq:beta}
\end{equation}
A positive value indicates an overall response along the reference differences;
it does not require agreement for each artist pair. Values above one indicate
exaggeration along that direction. To measure error, multiply the discrepancies
from the two repeats rather than squaring either repeat's discrepancy:
\begin{equation}
\begin{gathered}
\widehat D_{ms}=\frac1H\sum_a
\langle d_{msa1}-r_a,\ d_{msa2}-r_a\rangle,\\
\widehat D_m=\frac1S\sum_s\widehat D_{ms}.
\end{gathered}
\label{eq:distortion}
\end{equation}
Under stable conditional means and independent, mean-zero repeat errors,
noise that does not recur across requests contributes zero in expectation.
The expected score is then the normalized squared error of the scene-conditional
mean contrasts. Individual estimates can be negative; they are retained, not
interpreted as better-than-perfect recovery. The $D=1$ no-contrast benchmark
is a theoretical predictor, not the observed generic-painting arm.

Every scene is compared with the same pooled historical contrast. Scene-varying
artist responses therefore incur error even when their average matches that
target; this need not be an artistic mistake. We report both scene-wise error
and error after scene averaging (Appendix~\ref{app:score-details}).

\subsection{Model Comparisons and Uncertainty}
The prespecified 31-feature family contains six aligned responses and 15
model-pair error differences. Paired-scene Student intervals use a Bonferroni
adjustment for all 21, targeting approximate 95\% simultaneous coverage under
the stated model. Individual features, painters and outputs are not independent
replicates. Absolute-$D$ intervals are unadjusted descriptive summaries.

Inference conditions on the finite reference panel, scaling and 14 authored
briefs, assuming stable conditional means and independent scene/repeat error
blocks. Scene standard errors also include fixed scene heterogeneity and do
not isolate fresh-request variance. Two repeats cannot verify independence,
and the Student approximation guarantees neither actual service coverage nor
inference to unseen scenes. Under an assumed common trace correlation, the
Nano Banana 2--FLUX error point order crosses at $.251$
(Appendix~\ref{app:repeat-covariance}); actual dependence remains unidentified.
Declared scene, reference, feature-family and title-class sensitivities are
specified in Appendix~\ref{app:estimation}. They do not match compositions or
supply additional service observations.

\subsection{What the Additional Comparisons Test}\label{sec:diagnostics}
Artist-free and generic-painting arms characterize the removed common response.
Retrospective diagnostics examine all painter pairs, scenes and reference views.
A magnitude check fits one nonnegative multiplier of a model's centered
contrasts on 13 scenes and scores the omitted scene. It changes feature
contrasts, not images or references. These diagnostics do not expand the
confirmatory family (Appendix~\ref{app:diagnostics}).

\section{Results}\label{sec:results}
\subsection{A Response in the Right Direction Can Still Miss the Target}
All 1,008 outputs yield valid measurements. Every configuration's aligned-response
interval lies above zero under the stated inference model (Appendix
Table~\ref{tab:specificity-models}), but positive alignment need not give low error.
GPT Image 2 nearly matches the reference strength,
$\beta=.999\;[.893,1.104]$, while total error is $D=1.226$.
FLUX has lower alignment ($.470$) and the lowest error estimate ($.801$).
Only Flare--FLUX and Sunburst--FLUX resolve positive error differences after
the 21-comparison adjustment; the other 13 intervals include zero.
No configuration is established to beat the no-contrast error benchmark of one.
These are full-frame comparisons; source correction changes which differences
are resolved (Section~\ref{sec:source-correction}).

\subsection{Most Scene-averaged Change Is Common to the Painter Names}\label{sec:shared-change}
After averaging scenes, named-minus-artist-free shifts comprise a four-name mean
and painter-specific departures. The mean includes the oil-painting clause and
any common naming response: 82.5--95.7\% of repeat-corrected squared change
(Table~\ref{tab:shared-change}). Retrospectively changing the baseline to generic
oil painting isolates the additional naming change: 66.7--88.4\% is shared.
The painter-specific component is unchanged. Deleting each scene and recomputing
the ratios gives 65.9--90.2\% across models; these are sensitivity ranges,
not confidence intervals (Appendix~\ref{app:direct-naming}).

\begin{table*}[htbp]
\caption{Scene-averaged, repeat-corrected change. Shared is the four-name mean shift; between-name is departure from it. The artist-free baseline includes the painting clause; the retrospective generic-baseline row isolates additional naming. The last two rows compare the shared named-minus-free shift with generic-minus-free: corrected cosine measures directional agreement, and the ratio divides their corrected squared magnitudes, shared by generic.}
\label{tab:shared-change}
\centering\small\setlength{\tabcolsep}{3pt}
\begin{tabularx}{\linewidth}{@{}l*{6}{>{\centering\arraybackslash}X}@{}}
\toprule
& \shortstack{GPT Image\\1} & \shortstack{GPT Image\\2} & Flare & Sunburst & \shortstack{Nano\\Banana 2} & \shortstack{FLUX.2\\Max}\\\midrule
Free baseline: shared (\%) & 82.5 & 85.8 & 84.8 & 83.9 & 88.1 & 95.7\\
Free baseline: between-name (\%) & 17.5 & 14.2 & 15.2 & 16.1 & 11.9 & 4.3\\
Generic baseline: shared (\%) & 71.3 & 70.9 & 71.1 & 66.7 & 69.1 & 88.4\\
\midrule
Corrected cosine & 0.689 & 0.774 & 0.766 & 0.837 & 0.930 & 0.916\\
Squared-size ratio & 1.986 & 1.960 & 2.816 & 3.356 & 3.692 & 4.138\\
\bottomrule\end{tabularx}

\end{table*}


The shared named-minus-free shift points broadly along the generic oil-painting shift
(noise-corrected cosine $.689$--$.930$; 1 means parallel), but its repeat-corrected
squared magnitude is 1.96--4.14 times as large. The generic and additional naming
shifts interact: their signed cross term is negative for GPT Image 1 and positive
for the other five configurations (Appendix~\ref{app:direct-naming}). Their
squared magnitudes therefore do not add as separate contributions.
Computing the free-baseline decomposition within scenes instead gives
88.6--97.2\% shared change (Appendix~\ref{app:estimation}).


\subsection{A Separate Collection Delimits the Common-majority Result}\label{sec:cross-cohort}
We apply a fixed retrospective decomposition to all 2,000 retained SD-Turbo
outputs from an earlier collection: 25 repeat blocks, 16 scenes and the four
names plus a baseline already requesting oil painting. Naming inserts only
``by [painter]'' into a different template. Within each scene/block, all five
conditions share a seed; cross-products therefore use distinct complete
blocks, allowing correlated arms within a block. For block vectors $v_b$,
$U(v)=\sum_{b\ne b'}v_b^\top v_{b'}/[25(24)]$ replaces the two-repeat product.
The model, resolution and generation mechanism also differ from the main
collection (Appendix~\ref{app:cross-cohort}).

After scene averaging, the common component is 64.2\% of naming change,
with 63.9--64.6\% across all 25 block deletions. It remains a majority in
color (57.5\%) and spatial coordinates (84.9\%), but not texture (36.6\%;
deletion range 36.2--36.9\%). Computing within scenes gives 72.4\% overall
and 48.5\% for texture. Thus common-majority change persists in this separate
collection's complete metric, but does not hold for every feature family.

SD-Turbo's overall alignment is $.618$; pooled error is $.917$ and scene-wise
error $1.637$. Monet--Sisley alignment is $.471$, unlike the weak hand-feature
alignment in five main configurations. These are descriptive contrasts,
not checkpoint rankings. The image hashes are disjoint from the main collection,
but these pixels had prior distance analyses and reuse the same historical
target. This is a retrospective cross-collection check, not prospective
confirmation, a test of closed-service independence, or a shared-family control.

\subsection{Uneven Painter Distinctions in the 31 Features}
Monet--Sisley alignment ranges from $-.249$ to $.129$ in five configurations,
versus $.402$ for GPT Image 2 (Appendix Figure~\ref{fig:artist-pairs}). Yet GPT
Image 1 has lower pair error ($.56$ versus $1.16$). Removing C\'ezanne lowers
FLUX's overall alignment from $.470$ to $.096$, and GPT Image 2's from $.999$
to $.651$. Swapping Monet and Sisley lowers original-target error for Flare,
Sunburst and Nano Banana 2. Reversing a negatively aligned pair lowers its
error without changing its size; source correction can change individual signs.

\subsection{Scene Averaging and Contrast Magnitude Change the Ordering}
The error target matters even within the same coordinates.
Averaging scenes gives Nano Banana 2 lower error than FLUX ($.723$ versus $.761$),
whereas scene-wise scoring reverses that order because Nano Banana 2 varies more
across scenes. Rescaling centered contrasts by a multiplier fitted on the other
13 scenes instead gives GPT Image 2 the lowest held-scene error ($.554$ versus
FLUX's $.714$). These operations change the evaluated feature contrasts, not
images; lower error is not evidence of improved fidelity. The full decomposition,
calibration and influence checks are in Appendix~\ref{app:extended-magnitude}.

\subsection{Which Comparisons Survive Reference Changes?}\label{sec:source-correction}
FLUX retains the lowest uncalibrated 31-feature scene-error point estimate
across scene deletions, windows, reference targets and feature weighting
(Appendices~\ref{app:declared-sensitivities} and~\ref{app:diagnostics}).
AI assistants audited all 870 historical reproductions and selected peripheral
artifact-removal crops for 90 primary and 41 development works, without human
verification. With corrected regions and refitted scaling, FLUX's error changes
from $.801$ to $.872$ and all six aligned-response intervals remain positive.
The Sunburst--FLUX error interval now includes zero; Flare--FLUX remains
positive, and GPT Image 1--FLUX becomes positive. The latter is a retrospective
sensitivity, not an added confirmatory finding
(Appendix Table~\ref{tab:source-comparison}).
Weak Monet--Sisley alignment persists, but individual signs change.
Replacing title-derived with visual content labels for 11 nonmixed briefs
also retains FLUX's lowest point estimate. These corrections address visible
source regions and coarse content, not reproduction calibration or matched
historical compositions.

\subsection{What Changes in Learned Representations?}\label{sec:learned-results}
We remeasure all original images using CLIP ViT-L/14 \citep{radford2021}
and the public CSD ViT-L style checkpoint \citep{somepalli2024style}.
This retrospective protocol was fixed before learned outcomes. Each encoder
uses its native 224-pixel center crop and 768-dimensional unit vectors, without
IQR scaling. This changes representation and native preprocessing together;
the hand-feature square-window sensitivity in Appendix~\ref{app:direct-naming}
does not isolate encoder effects. Reference prototypes are means of the same 649 works;
Appendix~\ref{app:learned-audit} specifies checkpoints, checks and sensitivities.
CSD's release warns of a discrepancy from published benchmark results, so
we evaluate this released checkpoint without claiming benchmark reproduction.

\emph{Common movement also dominates prototype gain.}
Let $g_a$ and $\mu_a$ be mean named and reference embeddings, $g_0$ the generic
mean, and bars denote equal-painter averages. Using unnormalized prototypes,
mean own-painter cosine gain decomposes exactly as
\begin{equation}
\frac14\sum_a g_a^\top\mu_a-g_0^\top\bar\mu
=(\bar g-g_0)^\top\bar\mu+\frac{H\widehat\beta}{4}.
\label{eq:prototype-main}
\end{equation}
The first term is common movement; the second is labeled agreement.
Across configurations, the common term supplies 73.2--83.8\% of CLIP gain and
54.2--79.7\% of CSD gain (Table~\ref{tab:learned-main}). These are contributions
to cosine gain, distinct from the squared-change shares in
Table~\ref{tab:shared-change}. The latter are also majority-common:
76.1--82.3\% in CLIP and 68.8--77.4\% in CSD for naming beyond generic painting.

% Generated by paper/make_icml_learned_tables.py; do not edit numerical cells.
% Descriptive same-image audit; no cross-representation perceptual ranking.
% analysis_sha256: 06c77de4be133afc203dda122b2816d4f866e1389999f18bc94b3d9c668f28f7
% hand_review_v3_sha256: 79fa5ca28900e852c072e2a52d35f3cb400ee053935d857729b4ad4aa3bc542c
% input_sha256 reports/painter_learned_audit_v1/embeddings_clip.npz: d6e3d49bde6ca7d75fb3603cf7f3ecf7d46ea979b21a3b1e75435a504acd9945
% input_sha256 reports/painter_learned_audit_v1/embeddings_csd.npz: 65f428ed866916dca12ea0bc8148f1e0b9122343d19181e56f03543b5506085b
% input_sha256 reports/painter_learned_audit_v1/extraction_clip.json: 30a2aa7f93fcaa2e23413837ba3f75cf79b7ca3fca3f450a7b5ff4cddb7a22d5
% input_sha256 reports/painter_learned_audit_v1/extraction_csd.json: 5a55e4337b9588ef45bce46037dbe7ed123f89113d92da1e8d1c738eee0d6945
% input_sha256 reports/painter_learned_audit_v1/inputs.json: bd9cb00279fffb741b3579e3234ad034bb593884925ea40893c717612eccca57
% input_sha256 src/latent_art_bench/painter_learned_analysis_v1.py: d4b083896d7439a7018aa474969d588d0aaaf6eb170bfd26d477d1d31eb9eb46
% input_sha256 src/latent_art_bench/painter_learned_audit_v1.py: 156750d6d7099fb0a87b6ca789f5892d555ebca04fc58195aec85a8ea2a1c0c8
% input_sha256 studies/painter_learned_audit_v1/PLAN.md: c4af957b1922853b9541c02041d3a1b887a2f660dff709d3efb70710025766c0
\begin{table*}[t]
\caption{Same-image descriptive comparison using 649 primary references. Hand-31 uses full-frame features; CLIP and CSD use their native center crops of the original sources. Pair $\widehat\beta=1$ matches each representation's Monet--Sisley reference amplitude. Common prototype gain is the common contribution divided by named-minus-generic diagonal prototype gain (Eq.~\ref{eq:learned-prototype-gain}), not a squared-change share. Accuracy uses unit-normalized prototypes to classify the 112 named images per configuration (four painters; 25\% chance under uniform guessing). All cells are descriptive points; representations do not share a perceptual scale.}
\label{tab:learned-main}
\centering\small\setlength{\tabcolsep}{4pt}
\begin{tabularx}{\textwidth}{@{}X*{7}{>{\centering\arraybackslash}p{.081\textwidth}}@{}}
\toprule
& \multicolumn{3}{c}{Monet--Sisley $\widehat\beta$} & \multicolumn{2}{c}{\shortstack{Common prototype\\gain (\%)}} & \multicolumn{2}{c}{\shortstack{Macro accuracy\\(\%)}}\\
\cmidrule(lr){2-4}\cmidrule(lr){5-6}\cmidrule(l){7-8}
Requested configuration & Hand-31 & CLIP & CSD & CLIP & CSD & CLIP & CSD\\
\midrule
GPT Image 1 & $0.074$ & $0.301$ & $0.446$ & $74.5$ & $72.7$ & $60.7$ & $72.3$\\
GPT Image 2 & $0.402$ & $0.532$ & $0.587$ & $73.2$ & $54.2$ & $68.8$ & $75.9$\\
GPT Image 2.5 Flare & $-0.011$ & $0.339$ & $0.451$ & $78.8$ & $65.0$ & $59.8$ & $62.5$\\
GPT Image 2.5 Sunburst & $-0.249$ & $0.385$ & $0.323$ & $77.7$ & $67.4$ & $58.9$ & $61.6$\\
Nano Banana 2 & $-0.044$ & $0.318$ & $0.400$ & $83.8$ & $79.7$ & $51.8$ & $50.9$\\
FLUX.2 Max & $0.129$ & $0.299$ & $0.414$ & $81.7$ & $77.6$ & $41.1$ & $39.3$\\
\bottomrule\end{tabularx}
\end{table*}


\emph{Monet--Sisley alignment depends on representation.}
Monet--Sisley alignment is positive for all six configurations in both learned
representations: $.299$--$.532$ in CLIP and $.323$--$.587$ in CSD. Thus the
weak or reversed response measured by the 31 features cannot establish a
representation-independent failure to distinguish these painters. CLIP gives
GPT Image 1 the lowest scene-error estimate ($.847$); CSD gives FLUX the
lowest ($.714$). Distances across representations do not share a perceptual scale.

\emph{Proximity and recognition answer different questions.}
Classifying each named output by its closest unit-normalized reference prototype
gives the accuracies in Table~\ref{tab:learned-main}. Nano Banana 2 has the largest
CLIP prototype gain ($.1121$) but 51.8\% accuracy, versus GPT Image 2's smaller
gain ($.0996$) and 68.8\% accuracy. In CSD, GPT Image 2 has the smallest gain
($.1659$) and highest accuracy (75.9\%). These descriptive reversals show why
unpartitioned proximity gain is insufficient for evaluating labeled response.

Classifying the 221 development works with these primary prototypes yields
79.8\% macro accuracy in CLIP and 79.6\% in CSD. Appendix~\ref{app:learned-audit}
reports their painter confusions, reference energies and all absolute
common/labeled gain components. These controls describe panel separability,
not generated-image fidelity.

The appendix reports every configuration under both representations, original
and audited regions, and primary and development targets. These related
encoders share a CLIP backbone and unknown training-image overlap; the same
images supply all comparisons. They add a representation sensitivity, not a
fresh service replication or perceptual validation. No new significance claims
are attached to this audit.

\subsection{Common Offsets Can Change Prompted-name Decisions}\label{sec:transfer}
A decision task tests a consequence beyond score disagreement: recovering the
known painter clause of a generated image. This retrospective extension was
chosen after inspecting baseline confusions; its rules were fixed before
computing corrected outcomes. It uses all configurations and both encoders.
For each held-out scene, pool the 104 named outputs from the other 13 scenes,
excluding both held-scene repeats. Their mean $\bar g_{-s}$ ignores individual
painter assignments within the known balanced pool. With raw real-art centroids
$\mu_a$ and normalized directions $u_a=\mu_a/\|\mu_a\|$, set
\begin{equation}
 t_{-s}=\tfrac14\sum_a\mu_a-\bar g_{-s},\qquad
 \hat a=\arg\max_a (x+t_{-s})^\top u_a.
 \label{eq:transfer-main}
\end{equation}
The translation scale is fixed at one. It aligns mixture means and changes
class intercepts by $t_{-s}^\top u_a$; it leaves centered painter differences
unchanged within each fold. This cross-domain offset is not the
named-minus-generic effect in Equation~\ref{eq:prototype-main}: free and generic
outputs never enter its fit.

% Generated by paper/make_icml_transfer_tables.py; retained outcomes only.
% analysis_sha256: e1d1fbb924feb797f67e0907677511ed7c745335a4f2827e99d2f286e6be5448
% inputs_sha256: bdcf38bb09da53c890a1a62ce268c0336824ccadfb6ba90ee711bdaf1a6f1baf
% generator_sha256: eaefcdefea1cc8681bcd4c093a317222df52b46da74ce138b75b04270f4a4481
% bound_input_sha256 pyproject.toml: 5eb878e4cd9607f6dde074f4a341d54889f6405d65621821cdc4beb6713a2bf8
% bound_input_sha256 reports/painter_learned_audit_v1/analysis.json: 06c77de4be133afc203dda122b2816d4f866e1389999f18bc94b3d9c668f28f7
% bound_input_sha256 reports/painter_learned_audit_v1/csd_adapter_inputs_v2.json: 29107601907b257798a24bc01615c86e5984e11de0843c3dc8b4eb7685cf2567
% bound_input_sha256 reports/painter_learned_audit_v1/embeddings_clip.npz: d6e3d49bde6ca7d75fb3603cf7f3ecf7d46ea979b21a3b1e75435a504acd9945
% bound_input_sha256 reports/painter_learned_audit_v1/embeddings_csd.npz: 65f428ed866916dca12ea0bc8148f1e0b9122343d19181e56f03543b5506085b
% bound_input_sha256 reports/painter_learned_audit_v1/extraction_clip.json: 30a2aa7f93fcaa2e23413837ba3f75cf79b7ca3fca3f450a7b5ff4cddb7a22d5
% bound_input_sha256 reports/painter_learned_audit_v1/extraction_csd.json: 5a55e4337b9588ef45bce46037dbe7ed123f89113d92da1e8d1c738eee0d6945
% bound_input_sha256 reports/painter_learned_audit_v1/inputs.json: bd9cb00279fffb741b3579e3234ad034bb593884925ea40893c717612eccca57
% bound_input_sha256 reports/painter_learned_audit_v1/validation_clip.json: f298280c15d0801c0918a85e1e82c889e03bda90249e1a1072e227064fcf049d
% bound_input_sha256 reports/painter_learned_audit_v1/validation_csd.json: 40024de26abeb12a40381f00392c6e94019c5ef056218cd4704079984b023978
% bound_input_sha256 src/latent_art_bench/painter_learned_analysis_v1.py: d4b083896d7439a7018aa474969d588d0aaaf6eb170bfd26d477d1d31eb9eb46
% bound_input_sha256 src/latent_art_bench/painter_learned_audit_v1.py: 156750d6d7099fb0a87b6ca789f5892d555ebca04fc58195aec85a8ea2a1c0c8
% bound_input_sha256 src/latent_art_bench/painter_learned_csd_v1.py: 787b95697b6c48910a880c66961af4ea21c05d77184ddd0e8e7d2ff27395fb3e
% bound_input_sha256 src/latent_art_bench/painter_prototype_transfer_v1.py: 4a8f78bf438615c86cdce687469904bc5c191c13187e40e5d12be3e4747bb2e2
% bound_input_sha256 studies/painter_learned_audit_v1/PLAN.md: c4af957b1922853b9541c02041d3a1b887a2f660dff709d3efb70710025766c0
% bound_input_sha256 studies/painter_prototype_transfer_v1/PLAN.md: c7f02a6d65910bc6c8e1a1e79bac0f6c16143d0ba4176e69dc4455d3bf89e459
% bound_input_sha256 studies/painter_prototype_transfer_v1/PLAN_pre_design_qa.md: d5b9a549bc960215473eb13b838907a9389c81f6a73b7d4a7268d815571d2aaa
% bound_input_sha256 tests/painter_prototype_transfer_v1/test_membership_protocol.py: 8145e5c258f27f7ecb33f338edc34c7ee531cab563d15eb8aba5752413ccadb1
% bound_input_sha256 tests/painter_prototype_transfer_v1/test_rules.py: 43fc8cc7c92f6b09484cfb5cebe064262f6341400ff5284d753c2c39f7047573
% bound_input_sha256 uv.lock: d78b12a1ae92d586b9945d4bf50ed4398300780c462867e634abaab78941eb5a
\begin{table}[t]
\caption{Held-scene prompt-name accuracy (\%), original sources and primary references. B: reference prototypes; T: fixed common translation; G: supervised generated prototypes. T and G train on the other 13 scenes, excluding both held-out repeats. G uses painter labels and has more information than T. Each configuration has 112 named images. GPT 2.5 abbreviates GPT Image 2.5. These retrospective results concern prompt-name recovery, not perceptual fidelity.}
\label{tab:transfer-main}
\centering\footnotesize\setlength{\tabcolsep}{2pt}
\begin{tabularx}{\linewidth}{@{}X*{6}{>{\centering\arraybackslash}p{.095\linewidth}}@{}}
\toprule
& \multicolumn{3}{c}{CLIP} & \multicolumn{3}{c}{CSD}\\
\cmidrule(lr){2-4}\cmidrule(l){5-7}
Configuration & B & T & G & B & T & G\\\midrule
GPT Image 1 & $60.7$ & $55.4$ & $80.4$ & $72.3$ & $73.2$ & $84.8$\\
GPT Image 2 & $68.8$ & $66.1$ & $75.0$ & $75.9$ & $76.8$ & $94.6$\\
GPT 2.5 Flare & $59.8$ & $59.8$ & $77.7$ & $62.5$ & $74.1$ & $92.9$\\
GPT 2.5 Sunburst & $58.9$ & $62.5$ & $69.6$ & $61.6$ & $70.5$ & $90.2$\\
Nano Banana 2 & $51.8$ & $59.8$ & $63.4$ & $50.9$ & $61.6$ & $75.0$\\
FLUX.2 Max & $41.1$ & $53.6$ & $71.4$ & $39.3$ & $66.1$ & $75.0$\\
\bottomrule\end{tabularx}

\end{table}


On the original-view primary target, translation changes accuracy by $-5.4$
to $+12.5$ percentage points in CLIP,
with two declines, one tie and three gains; all six CSD changes are positive,
from $+0.9$ to $+26.8$ points. Equal-configuration means are $+2.7$ and $+10.0$
points, respectively. These summaries retain newly introduced errors:
CSD FLUX corrects 39 decisions and harms nine of 112, while CLIP GPT Image 1
corrects one and harms seven. Thus a common offset can affect recognition
without changing painter contrasts, but mean matching is not uniformly useful.

A supervised context rule fits one centroid per prompted painter on those same
13 training scenes. Its accuracy is higher than both reference rules in all 12
combinations (Table~\ref{tab:transfer-main}); it uses artist assignments that
translation does not. This demonstrates recoverable prompt-label structure
within generated images, not agreement with historical painting style.
Appendix~\ref{app:transfer} reports all confusions, painter recalls, scene changes
and reference sensitivities. Folds overlap and reuse previously inspected
images; these are descriptive out-of-fold results, not independent validation.

\subsection{Low Accepted Error Can Conceal Missing Painters}\label{sec:selective}
A further retrospective test asks when to withhold a reference assignment.
The fixed rule calibrates each painter's nonconformity
$s_a(x)=\max_{b\ne a}x^\top u_b-x^\top u_a$ on the other historical partition.
With fixed $\alpha=.10$, its threshold is the
$\lceil .9(n_a+1)\rceil$-th ordered calibration score. Retain the original
top-1 prediction only if it is the sole class meeting its threshold
(Appendix~\ref{app:selective}). Historical-to-generated exchangeability is not
established, so this rule provides no generated-image coverage or error guarantee.
No generated labels fit the gate. A comparator keeps the same number of named
queries by their ordinary top-two prototype margin, with fixed ID tie-breaking.

% Generated by paper/make_icml_selective_tables.py; retained outcomes only.
% analysis_sha256: f725b32b91f38fa8afd85454b3a6f0c8eb7464f5337b18373ce8c8f53e39fa2c
% inputs_sha256: f9d6b94e19d8004463f7b0524dd72865e22c18b0786cf48387cb9257e0528826
% audit_sha256: e3f26efb7286ad592e5ec30a50dd937690c29852d690c94b3baf09f31ae22569
% generator_sha256: 36191527ea108614a42cd949af9f3cdfbc16b01df7216fd40191198936079ec1
\begin{table}[!htbp]
\caption{Reference-calibrated abstention, primary CSD/original/primary setting. Coverage is accepted queries out of 112 named images per configuration; B is unrestricted top-1 error, A is accepted-gate error, and M is accepted error for the ordinary-margin comparator at exactly the same coverage. All rates are percentages. The fixed joint criterion fails coverage and painter retention; zero accepted errors do not establish a reliable four-painter attribution rule.}
\label{tab:selective-main}
\centering\footnotesize\setlength{\tabcolsep}{3pt}
\begin{tabularx}{\linewidth}{@{}Xrrrr@{}}\toprule
Configuration & Coverage & B & A & M\\\midrule
GPT Image 1 & $26.8$ & $27.7$ & $0.0$ & $0.0$\\
GPT Image 2 & $33.0$ & $24.1$ & $2.7$ & $2.7$\\
GPT 2.5 Flare & $24.1$ & $37.5$ & $0.0$ & $0.0$\\
GPT 2.5 Sunburst & $26.8$ & $38.4$ & $0.0$ & $0.0$\\
Nano Banana 2 & $45.5$ & $49.1$ & $25.5$ & $33.3$\\
FLUX.2 Max & $57.1$ & $60.7$ & $53.1$ & $48.4$\\
\bottomrule\end{tabularx}
\end{table}


The primary CSD/original/primary setting reduces equal-configuration mean
error by 26.0 percentage points relative to unrestricted prediction, but by
only 0.5 points relative to coverage-matched margin filtering
(Table~\ref{tab:selective-main}). Coverage is 24.1--57.1\%.
GPT Image 1 has zero errors among 30 accepted queries, yet 28 are C\'ezanne
and none is Sisley. Flare and Sunburst retain no Monet or Sisley queries.
FLUX retains more than half its queries but still has 53.1\% accepted error,
versus 48.4\% for the matched-margin comparator.
Figure~\ref{fig:selective-coverage} shows the complete painter coverage.

\begin{figure}[t]
\centering
\includegraphics[width=\linewidth]{icml_selective_coverage.pdf}
\caption{Primary CSD gate: accepted images per prompted painter, out of 28
in every cell. Selection concentrates on C\'ezanne in several configurations;
the zero aggregate errors for GPT Image 1, Flare and Sunburst in
Table~\ref{tab:selective-main} coexist with missing painter conditions.
Counts describe prompt-name coverage, not painter fidelity.}
\label{fig:selective-coverage}
\end{figure}

The pre-outcome operational criterion required at least 50\% coverage in each
configuration, some accepted query from every painter, and positive mean error
reductions against both comparisons. The primary setting fails the first two
conditions. All eight encoder/view/target settings fail at least the coverage
condition; this 50\% requirement is a chosen operating point, not a universal
standard. CLIP's original-primary mean reductions are 12.1 and 1.7 points,
with minimum coverage 49.1\%. The complete continuous results and scene deletions
remain in Appendix~\ref{app:selective}.

The primary gate also assigns a name to 14.3--50.0\% of artist-free controls
and 17.9--75.0\% of generic-painting controls. These are assignments without a
requested name, not verified visual misattributions. The operational lesson
is to report per-painter coverage and a matched-coverage comparison before
using reduced accepted error as evidence for a useful attribution policy.
The test does not establish that abstention generally fails, and no replacement
rule is selected after these outcomes.

\section{Discussion}\label{sec:discussion}
Detecting a name response and matching painter differences are distinct achievements.
FLUX's lowest uncalibrated error in the 31-feature metric coexists with weak
Monet--Sisley alignment, while learned representations yield positive pair alignment.
Model comparisons should therefore report overall and pairwise scores together
for the stated reference target, distinguishing estimates from resolved comparisons.
Reports should also state whether they score each scene, the average pattern,
or rescaled contrasts. For prompted-name recognition, the two held-scene
controls distinguish sensitivity to a pooled domain offset from supervised
within-generator separability. The abstention test adds a second decision check: lower accepted error can
coexist with missing painters and little advantage over margin filtering.
Neither lower contrast error nor better name recovery alone establishes closer
painter resemblance.

The intervention identifies differences between prompt conditions, not their
internal causes. Training composition, guidance and shared transformations remain
possible explanations.

\subsection{Scope and Remaining Limitations}
The target is the pattern of differences in a finite digital reference
collection. It mixes subject choices and reproduction processes with properties
of the paintings. Coarse content classes do not match composition, viewpoint,
career period or actual generated scene adherence. The source audit addresses
visible peripheral regions and title/visual-class disagreements, but AI
judgments, uncertain boundaries and remaining capture differences limit that
correction. One reproduction per work cannot separate effects of illumination,
sharpening, conservation or digitization source.

The 31 coordinates describe digital color, spatial organization and texture;
neither they nor the learned encoders establish physical brushwork or perceived
painter resemblance. CSD's style-tag list includes all four painters, and exact
training-image overlap is unknown for both encoders. Numerical replay and
representation sensitivities do not validate perceptual judgments.
The main inference conditions on four related painters, 14 authored briefs,
two repeats and one clause template. The SD-Turbo check changes the template
and uses 25 matched-seed blocks, but retains the painter group and historical
target. Neither collection has a shared-family prompt control. A common response
could represent appropriate family-level style; its prevalence alone does not
establish a misleading or inferior response. Simulations probe stated noise
models, not actual service independence. The image panels
make the comparison inspectable, while incomplete public exact-pixel access
still limits independent measurement reproduction.

\section{Conclusion}
Shared change, labeled reference agreement and prompt-name recovery lead to
different evaluation conclusions. Common-majority naming change appears across
the complete metrics of two collections, but fails for SD-Turbo texture.
Monet--Sisley alignment depends on representation. Mean translation changes
name recovery without changing painter contrasts, while a reference-calibrated
abstention rule can lower accepted error by retaining a painter-skewed subset.
Reports should state the reference and aggregation target, partition proximity
gain, and accompany selective error with per-painter coverage and an
coverage-matched comparator. These measured cases support that evaluation
practice; they do not establish a general ranking of painter fidelity.
